\section{Measuring whether a model reuses a planted error}
\label{sec:setup}

This section defines the protocol: how a model produces the reasoning we perturb, how we plant a false step, how we separate an error that can be re-read from one that must be recomputed, and how we score what the model does next.

We use short straight-line programs over integer variables as the reasoning substrate. Each program is a sequence of assignments in which a variable is set to a constant or to the sum of earlier variables, ending in a printed output. A program of this form has one correct execution trace, and every intermediate value is fixed by running it. This gives an unambiguous ground truth for what each step should say and lets us grade a continuation mechanically, which a natural-language task would not.

The model generates the trace we perturb. We give it the program, have it produce the execution trace itself, and keep only the traces it solves correctly. Every perturbation is therefore applied to the model's own prior reasoning rather than to a trace we wrote, so the model is continuing what it just produced rather than judging an outside artifact.

We replace exactly one line of the correct trace with a false value and leave the rest unchanged; the planted line keeps its variable and position and differs only in the value it asserts. We distinguish two kinds of planted error by what it would take to catch them. A \emph{readable} error contradicts a value written verbatim earlier in the trace, so it can be caught by comparison alone. A \emph{computed} error replaces the result of an operation, so catching it requires redoing that operation from its inputs. The inputs are present earlier in the trace in both cases; the only difference is whether the truth can be read off or must be recomputed.

Every planted value is audited to be genuinely false. A synthetic task makes it easy to plant a value that is wrong by construction yet happens to equal another true value in the trace, which would make absorption ambiguous, so we reject any perturbation whose planted value coincides with a re-readable true value and verify by re-execution that the planted line is false under the program.

We grade a continuation by re-executing the program and comparing the model's later values against both the true trace and the trace implied by the planted value. A continuation is \emph{absorbed} if its downstream values follow from the planted value, \emph{silently corrected} if they follow from the true value with no comment, and \emph{flagged} if the model explicitly notes the error. The three outcomes are decided by arithmetic, so no human or model judge enters the measurement. Absorption is the quantity of interest; silent correction and flagging are the two ways a model avoids it.

We elicit the continuation by having the model resume from the perturbed trace as its own partial output. For open-weight models, and for API models that allow it, we use prefilling, which places the perturbed trace in the assistant turn so the model continues token by token from the planted line. Some API models do not permit a prefilled assistant turn; for these we place the perturbed trace in the user turn and ask the model to continue it. On models that support both, the two methods give the same absorption on computed errors, so the presentation does not drive the result. One model is served only through a legacy completion endpoint, which is continuation by construction.

Figure~\ref{fig:protocol} shows the protocol on one program. The model has computed \texttt{u = 48} earlier in its own trace; a readable error inserts \texttt{u = 84} on the next line, and the model ignores it and continues from 48. A computed error instead changes a later step from its correct \texttt{j = 107} to \texttt{j = 97}, with the inputs \texttt{k = 70} and \texttt{c = 37} standing two lines above; the model takes 97 and carries it through every remaining step. We run this protocol on thirteen models from five developers, spanning a 1.5-billion-parameter open-weight model to current frontier systems.
